% -*- coding: utf-8 -*-
% XeLaTeX 编译
\documentclass[a3paper, landscape, 11pt, twocolumn]{article}

% 引入宏包
\usepackage{fontspec}
\usepackage[UTF8]{ctex}
\usepackage{geometry}
\usepackage{listings}
\usepackage{xcolor}
\usepackage{enumitem}
\usepackage{tabularx}
\usepackage{makecell}
\usepackage{fancyhdr}
\usepackage{amssymb}
\usepackage{lastpage}

% 设置页面边距 (A3横向双栏)
\geometry{left=2cm, right=2cm, top=2.5cm, bottom=2.5cm, columnsep=1.8cm}

% 颜色定义
\definecolor{codebg}{RGB}{245,245,245}
\definecolor{keywordcolor}{RGB}{0,0,150}
\definecolor{commentcolor}{RGB}{0,128,0}
\definecolor{stringcolor}{RGB}{163,21,21}

% Java代码块样式设置
\lstdefinestyle{JavaExam}{
    language=Java,
    basicstyle=\ttfamily\small,
    backgroundcolor=\color{codebg},
    keywordstyle=\color{keywordcolor}\bfseries,
    commentstyle=\color{commentcolor},
    stringstyle=\color{stringcolor},
    showstringspaces=false,
    breaklines=true,
    frame=single,
    rulecolor=\color{black!50},
    tabsize=4,
    xleftmargin=0.5em,
    xrightmargin=0.5em,
    aboveskip=0.5em,
    belowskip=0.5em,
    numbers=none
}

% 列表样式
\newlist{options}{enumerate}{1}
\setlist[options]{label=\Alph*., nosep, leftmargin=*, itemsep=0.3em}

% 页眉页脚设置
\pagestyle{fancy}
\fancyhf{}
\fancyhead[C]{\Large\bfseries 《Java程序设计》期末考试试卷}
\fancyfoot[C]{第 \thepage\ 页 共 \pageref{LastPage}\ 页}
\fancyfoot[R]{得分：\underline{\hspace{2cm}}}
\fancyfoot[L]{题号：\underline{\hspace{1cm}}}
\fancyheadoffset{0cm}
\fancyfootoffset{0cm}
\setlength{\headheight}{20pt}
\setlength{\footskip}{15pt}
\setlength{\headsep}{10pt}
\addtolength{\textheight}{-25pt}

% 重新定义页眉线
%\renewcommand{\headrulewidth}{0.4pt}

% 选择题选项命令
\newcommand{\chooseoption}[4]{\begin{options}\item #1\item #2\item #3\item #4\end{options}}

% 判断题
\newif\ifanswer
\answerfalse  % 考试模式，答案不显示
%\answertrue   % 阅卷模式，显示答案

% 填空题横线
\newcommand{\blank}{\underline{\hspace{2cm}}}

\begin{document}

% 试卷头部
\twocolumn[{
			\begin{center}
				\vspace*{-1cm}
				{\huge\bfseries 《Java程序设计》期末考试试卷（A 卷）}\\[0.8em]
				{\large 考试时间：120 分钟 \hspace{2em} 总分：150 分}\\[1em]
				\hrule
				\vspace{0.8em}
				\begin{tabularx}{0.95\linewidth}{|c|X|c|X|c|X|}
					\hline
					\makecell{题号} & \makecell{一                                                                        \\(40分)} & \makecell{二\\(20分)} & \makecell{三\\(20分)} & \makecell{四\\(30分)} & \makecell{五\\(40分)} \\
					\hline
					得分            & \hspace{1.5cm} & \hspace{1.5cm} & \hspace{1.5cm} & \hspace{1.5cm} & \hspace{1.5cm} \\
					\hline
				\end{tabularx}\\[1em]
				{\large 姓名：\underline{\hspace{3cm}} \hspace{1.5em} 学号：\underline{\hspace{4cm}} \hspace{1.5em} 班级：\underline{\hspace{3cm}}}
				\vspace{0.8em}
				\hrule
				\vspace{1em}
			\end{center}
		}]

\section*{一、单项选择题（共 20 小题，每小题 2 分，共 40 分）}
\textit{（在每小题给出的四个选项中，只有一项是符合题目要求的）}

\begin{enumerate}[label=\arabic*., itemsep=0.8em]
	\item 下列哪一项不是 Java 语言的特点？
	      \chooseoption{面向对象}{多继承}{平台无关性}{自动垃圾回收}

	\item 在 Java 中，源文件编译后生成的字节码文件扩展名是？
	      \chooseoption{.java}{.class}{.exe}{.jar}

	\item 下列关于 main 方法签名的定义，正确的是？
	      \chooseoption{public static int main(String[] args)}{public void main(String[] args)}{public static void main(String[] args)}{static public void main(String args)}

	\item 下列哪个关键字用于定义常量？
	      \chooseoption{static}{final}{const}{abstract}

	\item 表达式 \lstinline[style=JavaExam]|5/2| 的结果是？
	      \chooseoption{2.5}{2}{2.0}{编译错误}

	\item 关于数组的定义，错误的是？
	      \chooseoption{int[] a = new int[5];}{int a[] = {1,2,3};}{int a[5] = {1,2,3,4,5};}{int a[] = new int[]{1,2,3};}

	\item 在 Java 中，所有类的根类是？
	      \chooseoption{Object}{Class}{System}{Root}

	\item 关于构造方法的描述，错误的是？
	      \chooseoption{构造方法名必须与类名相同}{构造方法无返回值，但可以使用 void 声明}{构造方法可以重载}{构造方法在创建对象时自动调用}

	\item 下列哪个关键字用于实现接口？
	      \chooseoption{extends}{implements}{interface}{abstract}

	\item 关于抽象类的说法，正确的是？
	      \chooseoption{抽象类不能包含非抽象方法}{抽象类可以被实例化}{抽象类中的抽象方法必须在子类中实现（除非子类也是抽象类）}{抽象类中不能定义构造方法}

	\item 下列哪个异常是运行时异常（不受检查异常）？
	      \chooseoption{IOException}{SQLException}{NullPointerException}{ClassNotFoundException}

	\item 关于 finally 块的说法，正确的是？
	      \chooseoption{只有在 try 块没有异常时才会执行}{只有在 catch 块执行后才会执行}{无论是否发生异常，finally 块都会执行（除非 JVM 退出）}{可以没有 try 块单独存在}

	\item 下列哪个集合类是基于哈希表实现的，且不允许重复元素？
	      \chooseoption{ArrayList}{LinkedList}{HashSet}{TreeSet}

	\item 关于泛型的描述，错误的是？
	      \chooseoption{泛型可以提高类型安全}{泛型在运行时会被类型擦除}{泛型可以用于基本数据类型}{泛型方法可以在非泛型类中定义}

	\item 在 Java 中，启动线程的正确方法是？
	      \chooseoption{调用 run() 方法}{调用 start() 方法}{调用 begin() 方法}{调用 execute() 方法}

	\item 下列哪个类用于实现文件的字节输入流？
	      \chooseoption{FileReader}{FileInputStream}{BufferedReader}{FileWriter}

	\item 关于反射机制，下列说法正确的是？
	      \chooseoption{反射可以在运行时获取类的私有成员}{反射会提高程序性能}{反射只能在编译时使用}{反射不能动态创建对象}

	\item 在 JDBC 中，用于执行 SQL 查询并返回结果集的接口是？
	      \chooseoption{Statement}{PreparedStatement}{ResultSet}{Connection}

	\item 下列哪个布局管理器将容器划分为东、南、西、北、中五个区域？
	      \chooseoption{FlowLayout}{BorderLayout}{GridLayout}{CardLayout}

	\item 在 UDP 通信中，用于发送和接收数据报的类是？
	      \chooseoption{Socket}{ServerSocket}{DatagramSocket}{DatagramPacket}
\end{enumerate}

\section*{二、判断题（共 10 小题，每小题 2 分，共 20 分）}
\textit{（正确的选 A，错误的选 B）}

\begin{enumerate}[label=\arabic*., itemsep=0.6em]
	\item Java 程序必须显式导入 java.lang 包才能使用 System 类。 \underline{\hspace{1cm}}（A/B）
	\item 一个 Java 源文件中可以定义多个 public 类。 \underline{\hspace{1cm}}（A/B）
	\item 在方法重载中，返回值类型必须相同。 \underline{\hspace{1cm}}（A/B）
	\item 使用 final 修饰的方法可以被继承但不能被重写。 \underline{\hspace{1cm}}（A/B）
	\item 接口中的成员变量默认是 public static final 的。 \underline{\hspace{1cm}}（A/B）
	\item 捕获异常时，父类异常的 catch 块必须放在子类异常的后面。 \underline{\hspace{1cm}}（A/B）
	\item ArrayList 是线程安全的集合类。 \underline{\hspace{1cm}}（A/B）
	\item 泛型通配符 `?` 可以用于表示任意类型。 \underline{\hspace{1cm}}（A/B）
	\item 线程调用 sleep() 方法后不会释放持有的对象锁。 \underline{\hspace{1cm}}（A/B）
	\item 在 Swing 中，JFrame 默认的布局管理器是 FlowLayout。 \underline{\hspace{1cm}}（A/B）
\end{enumerate}

\section*{三、填空题（共 5 小题，每空 2 分，共 20 分）}

\begin{enumerate}[label=\arabic*., itemsep=0.8em]
	\item Java 的 8 种基本数据类型包括 byte、short、int、\blank、float、double、char 和 \blank。
	\item 面向对象三大特性是封装、\blank 和 \blank。
	\item 在 Java 中，使用关键字 \blank 导入包，使用 \blank 关键字声明包。
	\item 线程的生命周期包括新建、就绪、\blank、\blank 和死亡五种状态。
	\item 标准输入流对象是 System.\blank，标准输出流对象是 System.\blank。
\end{enumerate}

\section*{四、程序运行结果题（共 5 小题，每小题 6 分，共 30 分）}
\textit{（阅读下列程序，写出运行后的输出结果）}

\begin{enumerate}[label=\arabic*., itemsep=1em]
	\item
	      \begin{lstlisting}[style=JavaExam]
public class Test1 {
    public static void main(String[] args) {
        int sum = 0;
        for (int i = 1; i <= 5; i++) {
            if (i % 2 == 0) continue;
            sum += i;
        }
        System.out.println(sum);
    }
}
\end{lstlisting}
	      \textbf{输出：}\underline{\hspace{6cm}}

	\item
	      \begin{lstlisting}[style=JavaExam]
class Parent {
    String name = "Parent";
    public void show() {
        System.out.println(name);
    }
}
class Child extends Parent {
    String name = "Child";
    public void show() {
        System.out.println(name);
    }
}
public class Test2 {
    public static void main(String[] args) {
        Parent obj = new Child();
        obj.show();
    }
}
\end{lstlisting}
	      \textbf{输出：}\underline{\hspace{6cm}}

	\item
	      \begin{lstlisting}[style=JavaExam]
public class Test3 {
    public static void main(String[] args) {
        String s1 = "hello";
        String s2 = "hello";
        String s3 = new String("hello");
        System.out.println(s1 == s2);
        System.out.println(s1.equals(s3));
    }
}
\end{lstlisting}
	      \textbf{输出：}\underline{\hspace{6cm}}

	\item
	      \begin{lstlisting}[style=JavaExam]
public class Test4 {
    public static void main(String[] args) {
        try {
            int[] arr = {1, 2, 3};
            System.out.println(arr[3]);
        } catch (ArrayIndexOutOfBoundsException e) {
            System.out.println("数组越界");
        } finally {
            System.out.println("finally");
        }
        System.out.println("结束");
    }
}
\end{lstlisting}
	      \textbf{输出：}\underline{\hspace{6cm}}

	\item
	      \begin{lstlisting}[style=JavaExam]
class MyThread extends Thread {
    public void run() {
        for (int i = 0; i < 3; i++) {
            System.out.print(i + " ");
        }
    }
}
public class Test5 {
    public static void main(String[] args) {
        MyThread t1 = new MyThread();
        MyThread t2 = new MyThread();
        t1.start();
        t2.start();
    }
}
\end{lstlisting}
	      \textbf{输出可能的一种结果：}\underline{\hspace{6cm}}
\end{enumerate}

\newpage

\section*{五、编程题（共 2 小题，每小题 20 分，共 40 分）}
\textit{（要求代码完整、结构清晰、注释合理）}

\begin{enumerate}[label=\arabic*., itemsep=1em]
	\item \textbf{设计一个图形类层次结构} \\
	      编写一个抽象类 \texttt{Shape}，包含一个抽象方法 \texttt{double area()} 用于计算面积。从 \texttt{Shape} 派生出 \texttt{Circle} 类和 \texttt{Rectangle} 类。\\
	      \texttt{Circle} 类包含私有成员半径 \texttt{radius}，构造方法初始化半径，并实现 \texttt{area()} 方法。\\
	      \texttt{Rectangle} 类包含私有成员长 \texttt{length} 和宽 \texttt{width}，构造方法初始化，并实现 \texttt{area()} 方法。\\
	      在测试类 \texttt{Main} 中，创建一个 \texttt{Shape} 类型的数组，存入不同的图形对象，利用多态输出每个图形的面积。
	      \vspace{6cm}

	\item \textbf{学生成绩管理系统} \\
	      定义学生类 \texttt{Student}，属性包括学号（\texttt{id}）、姓名（\texttt{name}）、成绩（\texttt{score}）。要求：
	      \begin{itemize}
		      \item 提供无参和有参构造方法。
		      \item 为所有属性提供 getter 和 setter 方法。
		      \item 重写 \texttt{toString()} 方法，返回学生的完整信息。
		      \item 实现 \texttt{Comparable<Student>} 接口，按成绩降序排序。
	      \end{itemize}
	      在测试类中创建至少 5 个学生对象，存入数组，调用 \texttt{Arrays.sort()} 进行排序，并输出排序后的结果。
	      \vspace{6cm}
\end{enumerate}

\begin{center}
	\textbf{—— 试卷结束，祝考试顺利 ——}
\end{center}

\end{document}